% =====================================================================
%  Section 8.4: The worked example u_0 = -arctan(x) completed.
%  Uses the Lax--Oleinik formula (Remark 8.x) and the Cole--Hopf formula
%  (Section 8.2).  Optional, but Section 3.4 refers to it.
%  Code: code/arctan_example.py
% =====================================================================

\subsection{The worked example completed: $u_0=-\arctan x$}
\label{subsec:arctan-complete}

Sections~\ref{subsec:char-example} and \ref{subsec:shock-example} followed
$u_0(x)=-\arctan x$ up to its breakdown at $(x^*,t^*)=(0,1)$. The
Lax--Oleinik formula \eqref{eq:lax-oleinik} lets us continue the entropy
solution explicitly. The results of this subsection are worked
computations (with the proofs given in outline), not new theorems; the
numerical parts are produced by \texttt{code/arctan\_example.py}.

\subsubsection{Exact entropy solution after breakdown}

Here $V(y)=\int_0^y u_0=-\big(y\arctan y-\tfrac12\ln(1+y^2)\big)$, so the
critical points of $y\mapsto G(y)=V(y)+(x-y)^2/(2t)$ satisfy
\begin{equation}
  \label{eq:arctan-critical}
  G'(y)=-\arctan y+\frac{y-x}{t}=0
  \iff x=y-t\arctan y=\Phi_t(y),
\end{equation}
i.e.\ they are exactly the foot points $x_0=y$ of the characteristics
through $(x,t)$. For $t>1$ there are three of them when $|x|$ is small.

\begin{proposition}[Entropy solution for $u_0=-\arctan x$]
\label{prop:arctan-solution}
For $t<1$ the entropy solution is the classical solution
\eqref{eq:implicit-solution} (at $t=1$ it is continuous, with a vertical
tangent at $x=0$). For $t>1$ let $y^*(t)>0$ be the unique
positive root of $\arctan y=y/t$. Then the entropy solution is odd in $x$,
smooth for $x\ne0$, and for $x>0$
\begin{equation}
  \label{eq:arctan-outer}
  u(x,t)=-\arctan y,\qquad\text{where } y\ge y^*(t)\text{ is the unique solution of }
  y-t\arctan y=x,
\end{equation}
with a single \emph{stationary} shock at $x=0$ of strength
\begin{equation}
  \label{eq:arctan-strength}
  u_\ell(t)=-u_r(t)=\arctan y^*(t)\in(0,\pi/2).
\end{equation}
Moreover $u_\ell(t)=\sqrt{3(t-1)}\,(1+o(1))$ as $t\downarrow1$, and
$u_\ell(t)\to\pi/2$ as $t\to\infty$.
\end{proposition}

\begin{proof}[Proof (outline)]
For $t>1$, $G(\cdot)$ at $x=0$ has $G'(y)=y/t-\arctan y$, so $y=0$ is a
critical point with $G''(0)=(1-t)/t<0$ (a local \emph{maximum}). Since
$\arctan y/y$ decreases from $1$ to $0$ on $(0,\infty)$, there is exactly one
positive root $y^*$ of $\arctan y=y/t$, and by oddness of $G'$ the
minimisers at $x=0$ are the two points $\pm y^*$: a tie, hence a jump at
$x=0$, with $u_\ell=(0-(-y^*))/t=y^*/t=\arctan y^*$ and $u_r=-u_\ell$.
From $\arctan y>y/(1+y^2)$ for $y>0$ we get $1+y^{*2}>t$, so
$\Phi_t'(y)=1-t/(1+y^2)>0$ for $y\ge y^*$: on $[y^*,\infty)$ the map
$\Phi_t$ is strictly increasing from $0$ to $\infty$. For $x>0$, every
minimiser is $\ge y^*$ because the minimiser is nondecreasing in $x$
\cite[\S3.4.2]{evans2010}; so the minimiser is the unique root of
\eqref{eq:arctan-critical} in $[y^*,\infty)$, which gives
\eqref{eq:arctan-outer} (using $u=(x-y)/t=-\arctan y$), and in particular
there is no other jump. Oddness gives $x<0$. For the asymptotics, write
$\arctan y^*=y^*/t$ and $y^*\approx\sqrt{3(1-1/t)}$ from
$\arctan y=y-y^3/3+O(y^5)$ near $t=1$; and for $t\ge2$ the inequality $1+y^{*2}>t$ forces $y^*>1$, so
$y^*=t\arctan y^*\ge t\pi/4\to\infty$ and $\arctan y^*\to\pi/2$.
\end{proof}

Thus the shock is born at $(0,1)$ with zero strength, grows like
$\sqrt{3(t-1)}$, and saturates at $\pi/2$ --- the full jump between the
limits $u_0(-\infty)=\pi/2$ and $u_0(+\infty)=-\pi/2$ of the data, so
$u(\cdot,t)\to-\tfrac\pi2\operatorname{sgn}x$. This replaces the symmetry
remark of Section~\ref{subsec:shock-example} by an explicit shock strength
(Figure~\ref{fig:arctan}, left). Numerically, the parametrisation
\eqref{eq:arctan-outer} agrees with a brute-force minimisation of $G$ to
within the grid resolution ($\lesssim5\times10^{-4}$ for
$t\in\{0.5,1,2,5\}$), and $u_\ell/\sqrt{3(t-1)}=0.9999,\,0.9994,\,0.994,\,
0.945$ for $t-1=10^{-4},10^{-3},10^{-2},10^{-1}$.

\subsubsection{Vanishing viscosity for this example: a prediction and a check}

Theorem~\ref{thm:vanishing-viscosity} asserts that the Cole--Hopf solutions
$u^\nu$ converge to this entropy solution. Here we can predict the size of
the error, not just its vanishing.

\emph{Prediction.} Near the shock, $u^\nu$ should look like a stationary
viscous shock. By Proposition~\ref{prop:viscous-profile} (with $s=0$ and
$\Delta=2a$) such a shock is $-a\tanh(ax/2\nu)$, and its $L^1$ distance from the
sharp shock $-a\operatorname{sgn}x$ is $4\ln2\,\nu\approx2.77\,\nu$,
independent of the strength $a$. So if $u^\nu$ is locally a stationary shock of
slowly varying strength $a=u_\ell(t)$, the layer contributes
$\approx4\ln2\,\nu$ to $\|u^\nu-u\|_{L^1}$, and the rest of the error comes
from the smooth part, as an $O(\nu)$ viscous correction.

\emph{Check.} We evaluated \eqref{eq:cole-hopf-formula} with the above $V$
by quadrature (log-sum-exp stabilisation of the weights $e^{-G/2\nu}$) and
compared with Proposition~\ref{prop:arctan-solution} at $t=3$
(Figure~\ref{fig:arctan}, centre). Table~\ref{tab:arctan-visc} gives the
$L^1(-6,6)$ error: it is first order in $\nu$ and tends to $\approx3.11\,\nu$.
Splitting the integral at $|x|=0.05$ gives a layer contribution of
$2.78\,\nu$ (for $\nu=0.003,\,0.001,\,0.0003$) against the predicted
$2.77\,\nu$, within $0.5\%$, and a smooth-part contribution of
$\approx0.33\,\nu$.

\begin{table}[h]
\centering
\begin{tabular}{r r r r r r r r}
  $\nu$ & $0.1$ & $0.03$ & $0.01$ & $0.003$ & $0.001$ & $0.0003$ & $0.0001$ \\\hline
  $\|u^\nu-u\|_{L^1(-6,6)}$ & $3.16\!\times\!10^{-1}$ & $9.36\!\times\!10^{-2}$ &
    $3.11\!\times\!10^{-2}$ & $9.34\!\times\!10^{-3}$ & $3.11\!\times\!10^{-3}$ &
    $9.33\!\times\!10^{-4}$ & $3.10\!\times\!10^{-4}$ \\
  error$/\nu$ & 3.156 & 3.122 & 3.114 & 3.112 & 3.111 & 3.110 & 3.104 \\
  order (vs.\ previous $\nu$) & --- & 1.01 & 1.00 & 1.00 & 1.00 & 1.00 & 1.00
\end{tabular}
\caption{$L^1$ error of the Cole--Hopf solution against the entropy solution
at $t=3$.}
\label{tab:arctan-visc}
\end{table}

\begin{figure}[h]
\centering
\includegraphics[width=0.30\textwidth]{figures/arctan_shock_strength.pdf}
\includegraphics[width=0.66\textwidth]{figures/arctan_vanishing_viscosity.pdf}
\caption{Left: shock strength $u_\ell(t)=\arctan y^*(t)$ for
$u_0=-\arctan x$, with the small-$(t-1)$ asymptote and the limit $\pi/2$.
Centre: Cole--Hopf solutions at $t=3$ for decreasing $\nu$ against the
entropy solution (black). Right: $L^1$ error against $\nu$ (log--log).}
\label{fig:arctan}
\end{figure}

\emph{Is the layer error independent of the strength?} The prediction says
the layer contribution does not depend on $a$, whereas the data of
Table~\ref{tab:arctan-visc} all have $a=u_\ell(3)\approx1.32$. We repeated the
measurement at other times, i.e.\ other strengths
(\texttt{code/strength\_independence.py}, $\nu=3\times10^{-4}$), computing
the ratio of the error inside $|x|<\delta$ to $4\ln2\,\nu$ (prediction: $1$).

\begin{table}[h]
\centering
\begin{tabular}{c c c c c}
  $t$ & strength $a$ & outer slope $u_x(0^+)$ & ratio, $\delta=0.05$ & ratio, $\delta=0.005$ \\\hline
  1.1  & 0.518 & $-4.46$  & 1.464 & 1.087 \\
  1.5  & 0.967 & $-0.62$  & 1.054 & 1.006 \\
  3    & 1.324 & $-0.073$ & 1.004 & 1.0004 \\
  10   & 1.504 & $-0.005$ & 1.000 & 0.9999
\end{tabular}
\caption{Layer error divided by $4\ln2\,\nu$ for several shock strengths.}
\label{tab:strength-independence}
\end{table}

The ratio tends to $1$ as the window shrinks, at every strength tested. With
the wider window it deviates when the outer solution is steep beside the
shock (young shocks, $t$ close to $1$, large $|u_x(0^+)|$): then the smooth
part's $O(\nu)$ correction, which lies inside the window, is large, and at
$t=1.1$ the ratio at $\delta=0.005$ is still $9\%$ above $1$ and not
converged. So the prediction is confirmed where the shock is established
(its strength varies by a factor of $3$ across the table), and the early
deviations are explained by contamination from the outer solution, not by a
failure of the layer picture; this explanation is supported by the trend in
$\delta$ but not proved.

\noindent These are numerical observations for one data set,
not theorems; for general bounded data one only expects the weaker convergence
of Theorem~\ref{thm:vanishing-viscosity}. The prediction identifies the
layer part of the error exactly, but the $0.33\,\nu$ smooth-part constant is
only measured, not derived.